% GENERATED FILE. Edit canonical lessons, then run npm run generate:books.
% canonical-source-hash: fnv1a64:f542951ada0d9fb4
% canonical-lessons: JA-C55-fun, JA-C55-natsu, JA-C55-fuyu, JA-C55-itsumo, JA-C55-mada

\chapter{Minute, Summer, Winter, Always, Still}
\label{ch:ja-minute-summer-winter-always-still}
\hlchaptermodality{55}

\begin{chapteropening}
\textbf{By the end of this chapter:} \emph{I can say a minute, summer, winter, always and still, not yet.}

Complete the last lesson of chapter 55: I can say a minute, summer, winter, always and still, not yet.
\end{chapteropening}

\section[\texorpdfstring{\ja{ふん} (fun)}{ふん (fun)}]{\ja{ふん} (fun) — a minute}

\label{lesson:JA-C55-fun}



\begin{quote}
Before the new one: say the Japanese for every day, then the Japanese for a year; age.
\end{quote}

\subsection*{You'll want to know: \ja{ふん}}
\textbf{\ja{ふん}} — \emph{fun} — \textquotedblleft{}a minute\textquotedblright{}.

\begin{cousinweb}
\textbf{\ja{ふん}} (\emph{fun}) is Sino-Japanese (kango): the on'yomi of a character meaning \textquotedblleft{}divide; a minute\textquotedblright{}, a reading borrowed from Chinese (Middle Chinese \emph{pjun} in Baxter's transcription; Unihan, KANJIDIC2).
\end{cousinweb}

\begin{grammarlens}[title={What you've built}]
One more word for this chapter.
\end{grammarlens}

\subsection*{Your turn}
\begin{itemize}
\raggedright
  \item \emph{Say it:} \emph{fun}
  \item \emph{Say it:} \emph{fun}, once more
  \item \emph{Say it:} say \emph{fun}
  \item \emph{Recall:} say the Japanese for every day, then the Japanese for a year; age, then say \emph{fun} again
\end{itemize}

\begin{tcolorbox}[breakable,colback=teal!4,colframe=teal!35!black,title={Before you move on}]
What does \ja{ふん} mean? (\textquotedblleft{}A minute\textquotedblright{}.) Say it once more.
\end{tcolorbox}
\section[\texorpdfstring{\ja{なつ} (natsu)}{なつ (natsu)}]{\ja{なつ} (natsu) — summer}

\label{lesson:JA-C55-natsu}



\begin{quote}
Before the new one: say the Japanese for a year; age, then the Japanese for a minute.
\end{quote}

\subsection*{You'll want to know: \ja{なつ}}
\textbf{\ja{なつ}} — \emph{natsu} — \textquotedblleft{}summer\textquotedblright{}.

\begin{cousinweb}
\textbf{\ja{なつ}} (\emph{natsu}) is a native Japanese word, inherited rather than borrowed from Chinese, and its older origin is not settled. Its kanji was attached to it later as a kun'yomi, a native reading given to a Chinese character (JMdict, KANJIDIC2).
\end{cousinweb}

\begin{grammarlens}[title={What you've built}]
One more word for this chapter.
\end{grammarlens}

\subsection*{Your turn}
\begin{itemize}
\raggedright
  \item \emph{Say it:} \emph{natsu}
  \item \emph{Say it:} \emph{natsu}, once more
  \item \emph{Say it:} say \emph{natsu}
  \item \emph{Recall:} say the Japanese for a year; age, then the Japanese for a minute, then say \emph{natsu} again
\end{itemize}

\begin{tcolorbox}[breakable,colback=teal!4,colframe=teal!35!black,title={Before you move on}]
What does \ja{なつ} mean? (\textquotedblleft{}Summer\textquotedblright{}.) Say it once more.
\end{tcolorbox}
\section[\texorpdfstring{\ja{ふゆ} (fuyu)}{ふゆ (fuyu)}]{\ja{ふゆ} (fuyu) — winter}

\label{lesson:JA-C55-fuyu}



\begin{quote}
Before the new one: say the Japanese for a minute, then the Japanese for summer.
\end{quote}

\subsection*{You'll want to know: \ja{ふゆ}}
\textbf{\ja{ふゆ}} — \emph{fuyu} — \textquotedblleft{}winter\textquotedblright{}.

\begin{cousinweb}
\textbf{\ja{ふゆ}} (\emph{fuyu}) is a native Japanese word, inherited rather than borrowed from Chinese, and its older origin is not settled. Its kanji was attached to it later as a kun'yomi, a native reading given to a Chinese character (JMdict, KANJIDIC2).
\end{cousinweb}

\begin{grammarlens}[title={What you've built}]
One more word for this chapter.
\end{grammarlens}

\subsection*{Your turn}
\begin{itemize}
\raggedright
  \item \emph{Say it:} \emph{fuyu}
  \item \emph{Say it:} \emph{fuyu}, once more
  \item \emph{Say it:} say \emph{fuyu}
  \item \emph{Recall:} say the Japanese for a minute, then the Japanese for summer, then say \emph{fuyu} again
\end{itemize}

\begin{tcolorbox}[breakable,colback=teal!4,colframe=teal!35!black,title={Before you move on}]
What does \ja{ふゆ} mean? (\textquotedblleft{}Winter\textquotedblright{}.) Say it once more.
\end{tcolorbox}
\section[\texorpdfstring{\ja{いつも} (itsumo)}{いつも (itsumo)}]{\ja{いつも} (itsumo) — always}

\label{lesson:JA-C55-itsumo}



\begin{quote}
Before the new one: say the Japanese for summer, then the Japanese for winter.
\end{quote}

\subsection*{You'll want to know: \ja{いつも}}
\textbf{\ja{いつも}} — \emph{itsumo} — \textquotedblleft{}always\textquotedblright{}.

\begin{cousinweb}
\textbf{\ja{いつも}} (\emph{itsumo}) is native: \emph{itsu}, \textquotedblleft{}when\textquotedblright{}, with \emph{mo}, \textquotedblleft{}also, even\textquotedblright{}, so \textquotedblleft{}at any time\textquotedblright{} (JMdict). It is inherited rather than borrowed.
\end{cousinweb}

\begin{grammarlens}[title={What you've built}]
One more word for this chapter.
\end{grammarlens}

\subsection*{Your turn}
\begin{itemize}
\raggedright
  \item \emph{Say it:} \emph{itsumo}
  \item \emph{Say it:} \emph{itsumo}, once more
  \item \emph{Say it:} say \emph{itsumo}
  \item \emph{Recall:} say the Japanese for summer, then the Japanese for winter, then say \emph{itsumo} again
\end{itemize}

\begin{tcolorbox}[breakable,colback=teal!4,colframe=teal!35!black,title={Before you move on}]
What does \ja{いつも} mean? (\textquotedblleft{}Always\textquotedblright{}.) Say it once more.
\end{tcolorbox}
\section[\texorpdfstring{\ja{まだ} (mada)}{まだ (mada)}]{\ja{まだ} (mada) — still, not yet}

\label{lesson:JA-C55-mada}



\begin{quote}
Before the new one: say the Japanese for winter, then the Japanese for always.
\end{quote}

\subsection*{You'll want to know: \ja{まだ}}
\textbf{\ja{まだ}} — \emph{mada} — \textquotedblleft{}still, not yet\textquotedblright{}.

\begin{cousinweb}
\textbf{\ja{まだ}} (\emph{mada}) is a native Japanese word, inherited rather than borrowed from Chinese, and its older origin is not settled. It is usually written in kana; its kanji, when used, is a kun'yomi attached later (JMdict, KANJIDIC2).
\end{cousinweb}

\begin{grammarlens}[title={What you've built}]
One more word for this chapter.
\end{grammarlens}

\subsection*{Your turn}
\begin{itemize}
\raggedright
  \item \emph{Say it:} \emph{mada}
  \item \emph{Say it:} \emph{mada}, once more
  \item \emph{Say it:} say \emph{mada}
  \item \emph{Recall:} say the Japanese for winter, then the Japanese for always, then say \emph{mada} again
\end{itemize}

\begin{tcolorbox}[breakable,colback=teal!4,colframe=teal!35!black,title={Before you move on}]
What does \ja{まだ} mean? (\textquotedblleft{}Still, not yet\textquotedblright{}.) Say it once more.
\end{tcolorbox}
